\documentclass[12pt]{article}
\usepackage[utf8]{inputenc}
\usepackage{amsmath}
\usepackage{geometry}
\geometry{a4paper, margin=1in}

\title{Metodologías Ágiles: Scrum y Lean Six Sigma}
\author{}
\date{}

\begin{document}

\maketitle

\section*{Introducción}
Agile, Scrum y Lean Six Sigma son marcos de trabajo orientados a mejorar la eficiencia, calidad y adaptabilidad en proyectos. Mientras Agile establece una filosofía general, Scrum implementa un proceso iterativo y Lean Six Sigma combina eliminación de desperdicios con reducción de variabilidad.

\section*{Agile}
Agile es un conjunto de principios (Manifiesto Ágil, 2001) que prioriza:
\begin{itemize}
    \item Individuos e interacciones sobre procesos y herramientas.
    \item Software funcionando sobre documentación exhaustiva.
    \item Colaboración con el cliente sobre negociación contractual.
    \item Respuesta ante el cambio sobre seguir un plan.
\end{itemize}
Se utiliza en entornos dinámicos donde los requisitos evolucionan.

\section*{Scrum}
Scrum es el marco ágil más popular. Sus componentes clave:
\begin{itemize}
    \item \textbf{Roles:} Product Owner (prioriza el backlog), Scrum Master (facilita el proceso) y Development Team (autoorganizado).
    \item \textbf{Eventos:} Sprint (ciclo de 1–4 semanas), Daily Scrum, Sprint Review y Sprint Retrospective.
    \item \textbf{Artefactos:} Product Backlog, Sprint Backlog e Incremento.
\end{itemize}
Scrum entrega valor incrementalmente y permite inspección y adaptación continuas.

\section*{Lean Six Sigma}
Lean Six Sigma fusiona dos enfoques:
\begin{itemize}
    \item \textbf{Lean:} Elimina desperdicios (sobreproducción, esperas, defectos, etc.) y optimiza el flujo de valor.
    \item \textbf{Six Sigma:} Reduce la variación de procesos mediante DMAIC (Definir, Medir, Analizar, Mejorar, Controlar) y herramientas estadísticas.
\end{itemize}
Es ideal para procesos repetitivos de manufactura, logística o servicios transaccionales.

\section*{Comparación y uso combinado}
\begin{itemize}
    \item \textbf{Scrum} se adapta a proyectos de incertidumbre (software, I+D).
    \item \textbf{Lean Six Sigma} se aplica a procesos estables que requieren calidad y eficiencia.
    \item En la práctica, muchas organizaciones usan Scrum para desarrollo y Lean Six Sigma para operaciones, o integran herramientas de Lean (Kanban) dentro de Scrum.
\end{itemize}

\section*{Conclusión}
Agile y Scrum priorizan la adaptabilidad y entrega rápida; Lean Six Sigma se enfoca en la reducción de defectos y desperdicios. La elección depende del tipo de proyecto. Combinarlos puede potenciar la mejora continua y la satisfacción del cliente.

\end{document}
